% from ega0-15-fr.tex
\oldpage[0\textsubscript{IV-1}]{22}
Soit $u=(\psi,\theta):X\to Y$ un morphisme plat d'espaces annelés
($0_{\mathrm I}$, 6.7.1)~; soient $\mathscr{F}$ un
$\mathcal{O}_Y$-Module, $\mathscr{G}=u^*(\mathscr{F})$,
$(f_i)_{1\leq i\leq n}$ une suite d'éléments de
$\Gamma(Y,\mathcal{O}_Y)$, $(g_i)_{1\leq i\leq n}$ la suite de leurs images
$g_i=\theta(f_i)\in\Gamma(Y,u_*(\mathcal{O}_X))=\Gamma(X,\mathcal{O}_X)$. Si
la suite $(f_i)$ est $\mathscr{F}$-régulière, la suite $(g_i)$ est
$\mathscr{G}$-régulière. La réciproque est vraie si $u$ est un morphisme
fidèlement plat ($0_{\mathrm I}$, 6.7.8).
\end{proposition}

\begin{proof}
Soient $x$ un point de $X$, $y=u(x)=\psi(x)$~; on a
\[
  \mathscr{G}_x=\psi_x^*(\mathscr{F}_y)
  \otimes_{\psi_x^*(\mathcal{O}_y)}\mathcal{O}_x,
  \qquad\text{et}\qquad
  (g_i)_x=\theta_x^\#(\psi_x^*((f_i)_y))~;
\]
comme par hypothèse
$\theta_x^\#:\psi_x^*(\mathcal{O}_y)\to\mathcal{O}_x$ fait de
$\mathcal{O}_x$ un $\psi_x^*(\mathcal{O}_y)$-module plat, et que le foncteur
$\psi^*$ est exact, la première assertion résulte de l'assertion correspondante
de (15.1.14)~; il en est de même de la seconde, compte tenu de ce que
l'hypothèse que $u$ est fidèlement plat signifie que $\psi$ est surjectif et
que $\theta_x^\#$ fait de $\mathcal{O}_x$ un
$\psi_x^*(\mathcal{O}_y)$-module fidèlement plat.
\end{proof}

% (continued in ega0-16-fr.tex, before its first page marker)
\section{Dimension et profondeur dans les anneaux locaux noethériens}
\label{section:0.16-fr}

\subsection{Dimension d'un anneau}
\label{subsection:0.16.1-fr}

\begin{env}[16.1.1]
\label{0.16.1.1-fr}
On appelle dimension (ou dimension de Krull) d'un anneau $A$ et on note
$\dim A$ la dimension (combinatoire) de son spectre
$\operatorname{Spec}(A)$ (14.2.1)~; comme les parties fermées irréductibles de
$\operatorname{Spec}(A)$ sont les
$V(\mathfrak{p})$, où $\mathfrak{p}$ est un idéal premier de $A$ (Bourbaki,
\emph{Alg.~comm.}, chap.~II, \S~4, no~3, cor.~1 de la prop.~14), il revient au
même de dire que $\dim(A)$ est la borne supérieure des longueurs des chaînes
(14.1.1)
\[
  \mathfrak{p}_0\subset\mathfrak{p}_1\subset\cdots\subset\mathfrak{p}_n
\]
d'idéaux premiers de $A$. Si $A\neq0$, on a $\dim A\geq0$. Un anneau artinien
$\neq0$ (en particulier un corps) est de dimension 0 (Bourbaki,
\emph{Alg.~comm.}, chap.~IV, \S~2, no~5, prop.~7)~; un anneau de Dedekind qui
n'est pas un corps (en particulier $\mathbb{Z}$ ou un anneau $k[T]$ de
polynômes à une indéterminée sur un corps $k$) est de dimension 1 (Bourbaki,
\emph{Alg.~comm.}, chap.~VII, \S~2). Un anneau noethérien n'est pas
nécessairement de dimension finie [30].

Si $(\mathfrak{p}_\alpha)$ est la famille des idéaux premiers minimaux de $A$,
on a (14.1.2.1)
\[
\label{0.16.1.1.1-fr}
  \dim(A)=\sup_\alpha(\dim(A/\mathfrak{p}_\alpha)).
\tag{16.1.1.1}
\]
\end{env}

\begin{env}[16.1.2]
\label{0.16.1.2-fr}
Pour tout idéal $\mathfrak{a}$ de $A$, les idéaux premiers de
$A/\mathfrak{a}$ correspondent biunivoquement aux idéaux premiers de $A$
contenant $\mathfrak{a}$, donc
\[
\label{0.16.1.2.1-fr}
  \dim(A/\mathfrak{a})\leq\dim(A).
\tag{16.1.2.1}
\]
Si en outre $\mathfrak{a}$ n'est contenu dans aucun idéal premier minimal
$\mathfrak{p}_\alpha$ de $A$, toute chaîne d'idéaux premiers de $A$ contenant
$\mathfrak{a}$ peut être complétée par un des $\mathfrak{p}_\alpha$, donc on a
\[
\label{0.16.1.2.2-fr}
  1+\dim(A/\mathfrak{a})\leq\dim(A).
\tag{16.1.2.2}
\]
\end{env}

\begin{env}[16.1.3]
\label{0.16.1.3-fr}
